\subsubsection{Martini}
\noindent
\begin{minipage}{.49\textwidth}
    \noindent\textbf{\LARGE Martini\index{Martini}}
\end{minipage}
\hspace{.02\textwidth}
\begin{minipage}{.49\textwidth}
    \raggedleft
    \textbf{\large Gin}\\
    \sweettartguage{-0.1}\index[drinks]{Sweet!10\%!Martini}
\end{minipage}
\vspace{.5em}\hrule\vspace{.5em}
\noindent
\textbf{\large Flavor Profile:} Herbal\index[flavor-profile]{Herbal}, Juniper\index[flavor-profile]{Juniper}, Dry\index[flavor-profile]{Dry}, Hints of Citrus\index[flavor-profile]{Citrus}.
\vspace{.5em}\hrule\vspace{.5em}
\noindent
\textbf{\large History:} The Martini is the most popular cocktail in the world, and yet the perfect martini is wildly contested from one person to the next. The Martini is a variation of the Martinez that dates back to the 1860s. Classically, this cocktail would include more vermouth, 2 dashes of orange bitters, and always include gin. The modern version uses less vermouth, typically does not include orange bitters, and is garnished with a lemon twist or olives. So, experiment and find your perfect Martini!
\vspace{.5em}\hrule\vspace{.5em}
\noindent
\begin{minipage}[t]{.49\textwidth}
    \noindent
    \textbf{\large Ingredients}
    \begin{itemize}
        \item Gin 67,5ml\index[ingredient]{Gin}
        \item Dry Vermouth 22,5ml\index[ingredient]{Vermouth}
        \item Olive Juice (optional for a "dirty" Martini) 15ml\index[ingredient]{Olive Juice}
    \end{itemize}
    \noindent
    \textbf{\large Garnish}
    \begin{itemize}
        \item Lemon Twist\index[ingredient]{Lemon Twist}, Orange Twist\index[ingredient]{Orange Twist} or Pickled Onions\index[ingredient]{Pickled Onion}
    \end{itemize}
\end{minipage}
\hspace{.01\textwidth}\vrule\hspace{.01\textwidth}
\begin{minipage}[t]{.49\textwidth}
    \noindent
    \textbf{\large Tools/Glassware}
    \begin{itemize}
        \item Mixing Glass
        \item Hawthorne Strainer
        \item Bar Spoon
        \item Jigger
        \item Y-Peeler
        \item Coupe\index[glass]{Coupe} or Martini Glass\index[glass]{Martini Glass}
    \end{itemize}
\end{minipage}
\vspace{.5em}\hrule\vspace{.5em}
\noindent
\begin{minipage}{\textwidth}
    \noindent
    \textbf{\large Instructions}
    \begin{itemize}
        \item Place gin and vermouth into mixing glass.
        \item Add ice and stir for 30 seconds.
        \item Strain into coupe or martini glass.
        \item Finish with your choice of garnish.
    \end{itemize}
\end{minipage}
\vspace{.5em}\hrule\vspace{.5em}
\noindent
\textbf{\large Fun Tip:} A "dry" martini refers to using less vermouth, and a "wet" martini refers to adding more vermouth. Try them out and find your favorite combination!
\newpage
